\section{Reciprocal-lattice templates and \texorpdfstring{low-$\ell$}{low-l} alignment}
\label{app:moire_cmb}
Part~VI uses a reciprocal-lattice (``moir\'e'') template to motivate preferred-axis signatures in low-$\ell$ CMB multipoles. Here we record the elementary spherical-harmonic relation for sums of plane waves on the sky.

\begin{proposition}[Plane-wave template and spherical-harmonic coefficients]
\label{prop:plane_wave_harmonics}
Let $\hat{n}\in S^2$ denote a sky direction and consider a (real) scalar template of the form
\begin{equation}
    V(\hat n)=\sum_{k\in\mathcal{K}} V_k\,e^{i\,\vec{k}\cdot \hat n},\qquad V_{-k}=\overline{V_k},
\end{equation}
where $\mathcal{K}$ is a finite (or sufficiently summable) set of wavevectors. Expand $V$ in spherical harmonics, $V(\hat n)=\sum_{\ell,m} a_{\ell m}Y_{\ell m}(\hat n)$. Then
\begin{equation}
    a_{\ell m} = 4\pi\, i^{\ell}\sum_{k\in\mathcal{K}} V_k\, j_{\ell}(|k|)\,\overline{Y_{\ell m}(\hat{k})},
\end{equation}
where $j_{\ell}$ is the spherical Bessel function and $\hat{k}\equiv \vec{k}/|\vec{k}|$.
\end{proposition}
\begin{proof}[Reference]
This follows from the standard plane-wave expansion into spherical harmonics; see, e.g., \cite{ArfkenWeberHarris2013}.
\end{proof}

\begin{corollary}[Dominant axis implies aligned low-$\ell$ multipoles]
\label{cor:axis_alignment}
If the set of dominant wavevectors $\mathcal{K}$ is concentrated near a single axis $\hat{k}_\ast$ (for example, $\mathcal{K}$ is dominated by $\pm \vec{k}_\ast$), then the low-$\ell$ coefficients $a_{\ell m}$ are dominated by the corresponding spherical harmonics evaluated at $\hat{k}_\ast$. In particular, in the exactly axisymmetric case $V(\hat n)=f(\hat{k}_\ast\cdot \hat n)$, one can choose coordinates with $z$-axis along $\hat{k}_\ast$ so that $a_{\ell m}=0$ for all $m\neq 0$, implying that the quadrupole and octopole share the same preferred axis.
\end{corollary}
\begin{proof}[Reference]
This is the standard fact that axisymmetric functions on $S^2$ have only $m=0$ spherical-harmonic components in coordinates aligned with the symmetry axis; see, e.g., \cite{ArfkenWeberHarris2013}.
\end{proof}

